\subsection{等价对象的类}

设 \(\mathbb{R}\{x, y\}\) 是一个等价关系，它不必有图。显然，如果 \(x, x', y\) 是三个不同的字母，关系 \(\mathbb{R}\{x, x'\}\) 蕴含“ \(\mathbb{R}\{x, y\} \Longleftrightarrow \mathbb{R}\{x', y\}\) ”，因此也蕴含关系 \((\forall y)(\mathbb{R}\{x, y\} \Longleftrightarrow \mathbb{R}\{x', y\})\) 。根据方案S7（第一章，§5，第1号），如果我们令 \(\theta\{x\} = \tau_y(\mathbb{R}\{x, y\})\) ，关系 \(\mathbb{R}\{x, x'\}\) 蕴含

\[
\theta \{x \} = \theta \left\{x ^ {\prime} \right\}.
\]

注意，另一方面，根据定义 \(\mathbb{R}\{x, \theta\{x\}\}\) 是关系 \((\exists y)(\mathbb{R}\{x, y\})\) ，因此（第1号）等价于 \(\mathbb{R}\{x, x\}\) 。由此可知，关系 \((\mathbb{R}\{x, x\}\) 并且 \(\mathbb{R}\{x', x'\}\) 并且 \(\theta\{x\} = \theta\{x'\})\) 等价于 \(\mathbb{R}\{x, x'\}\) ，因为它通过S6（第一章，§5，第1号）蕴含关系

\[
(R \{x, x \} \text {   且   } R \{x ^ {\prime}, x ^ {\prime} \} \text {   且   } (R \{x ^ {\prime}, \theta \{x \} \} \iff R \{x ^ {\prime}, \theta \{x ^ {\prime} \} \})),
\]

进而蕴含 \((\mathbb{R}\{x, \theta\{x\}\}\) 和 \(\mathbb{R}\{x', \theta\{x\}\})\) ，最后通过传递性和对称性得到 \(\mathbb{R}\{x, x'\}\) 。并且由于反过来， \(\mathbb{R}\{x, x'\}\) 蕴含 \(\mathbb{R}\{x, x\}\) 和 \(\mathbb{R}\{x', x'\}\) ，该断言得证。术语 \(\theta\{x\}\) 被称为与 \(x\) 等价的对象类（关于关系 \(\mathbb{R}\) ）。现在假设 \(\mathbb{T}\) 是一个使得关系

\[
(\forall y) (R \{y, y \} \Longrightarrow (\exists x) (x \in T \text {   且   } R \{x, y \}))\tag{1}
\]

成立的项。那么关系 \((\exists x)(R\{x, x\}\) 并且 \(z = \theta\{x\})\) 关于 \(z\) 是集合化的。因为我们可以假设 \(x \in T\) 蕴含 \(R\{x, x\}\) ；只需将 \(T\) 替换为所有 \(x \in T\) （使得 \(R\{x, x\}\) ）组成的集合（注意到 \(R\{x, y\}\) 蕴含 \(R\{x, x\}\) ）。设 \(\Theta\) 是所有形如 \(\theta\{x\}\) （ \(x \in T\) ）的对象的集合 ( \(\S 1\) ，第6号）。假设 \(R\{y, y\}\) 成立；那么存在 \(x \in T\) 使得 \(R\{x, y\}\) 因此 \(\theta\{y\} = \theta\{x\} \in \Theta\) 。集合 \(\Theta\) 被称为关于 \(R\) 的等价对象的类的集合；并且对于每一个 \(x\) 使得 \(R\{x, x\}\) , \(\theta\{x\}\) 是唯一元素 \(z \in \Theta\) 使得 \(R\{x, z\}\) .

在同一假设下，设 \(\mathbf{A}\{x\}\) 是一个项，使得 \(\mathbf{R}\{x,y\}\) 蕴含 \(\mathbf{A}\{x\} = \mathbf{A}\{y\}\) 。则关系 \((\exists x)(\mathbf{R}\{x,x\}\) 并且 \(z = \mathbf{A}\{x\})\) 关于 \(z\) 是集合化的，因为 \(\mathbf{R}\{x,x\}\) （它等价于 \(\mathbf{R}\{x,\theta\{x\}\}\) ）蕴含 \(\mathbf{A}\{x\} = \mathbf{A}\{\theta\{x\}\}\) ，因此如果 \(\mathbf{E}\) 是形如 \(\mathbf{A}\{t\}\) （对于 \(t \in \Theta\) ）的对象组成的集合，那么 \(\mathbf{R}\{x,x\}\) 蕴含 \(\mathbf{A}\{x\} \in \mathbf{E}\) 。如果 \(f\) 是函数 \(t \to \mathbf{A}\{t\}\) ( \(t \in \Theta\) , \(\mathbf{A}\{t\} \in \mathbf{E}\) ），则关系 \(\mathbf{R}\{x,x\}\) 蕴含 \(\mathbf{A}\{x\} = f(\theta\{x\})\) .

特别地，如果 R 是集合 F 上的一个等价关系，我们可以取 A \(\{x\}\) 为 x 关于R（第2条）的等价类，且函数 f 随后是从 \(\Theta\) 到商集 F/R 上的双射；这证明了所引入术语的合理性。

* 例。设 \(\mathbb{R}\{x, y\}\) 为等价关系“\(x\) 并且 \(y\) 是 \(\mathbf{C}\)上两个具有相同有限维数的向量空间”；此关系没有图。当 \(\mathbf{T}\) 是 \(\mathbf{C}^{(\mathbf{N})}\) 的所有向量子空间的集合时，或当子集 \(\mathbf{T}'\) 的 \(\mathbf{T}\) 由空间 \(\mathbf{C}^n\) ( \(n \in \mathbf{N}\) ）组成时，它满足条件（1），并约定 \(\mathbf{C}^0\) 由 \(\mathbf{C}^{(\mathbf{N})}\) 以及 \(\mathbf{C}^n\) ( \(n > 0\) ) 是前 \(n\) 直和的组成部分 \(\mathbf{C}^{(\mathbf{N})}\) 采用这第二种选择，我们便有 \(\Theta = \mathbf{T}'_*\)

